% Figure-by-figure analysis of the V2V-BFT ablation studies.
% Every number here is traceable to benchmarks/ablations/results/ and every
% mechanism to a named file and line in src/ or bridge/.
\documentclass[11pt,a4paper]{article}
\usepackage[margin=2.5cm]{geometry}
\usepackage{booktabs}
\usepackage{graphicx}
\usepackage{hyperref}
\usepackage{xcolor}
\newcommand{\code}[1]{\texttt{\small #1}}
\newcommand{\caveat}[1]{\textcolor{red!60!black}{\textbf{#1}}}

\title{V2V--BFT Intersection Protocol:\\Figure-by-Figure Analysis}
\author{}
\date{27 August 2026}

\begin{document}
\maketitle

\section{Experimental setup}

Four approaches, \textbf{one lane each}, no turn pockets. Each vehicle draws its
exit uniformly from its approach's three valid exits via a SUMO
\code{routeDistribution}; the U-turn is absent from the distribution, so no
vehicle leaves the way it came. Roughly one third of vehicles therefore turn
left. \code{*.manager.seed} varies per repetition, so each repetition samples a
different turn assignment.

All vehicles depart at $t=0$: a single arrival burst, not continuous demand.
This is a consequence of the protocol completing \textbf{one consensus epoch}
per run --- \code{ResDBArrivalProtocol.cc} advances discovery for epoch~0 and a
long-lived multi-epoch scenario is not yet supported. Every claim below is
therefore bounded to single-round, burst-arrival conditions.

Repetition counts differ by study and are stated per figure. Min--max ranges are
not drawn on any figure; the spread is computed per cell and printed in
\code{figures/results.md}.

\section{Ablation 1 --- RSU units}

\textbf{9 repetitions} at $k{=}0$ (the condition all three panels plot);
3 repetitions across the $k$ sweep.

\subsection*{Observation}

Throughput and latency (stop to clearing) are \emph{indistinguishable} between the
with-RSU and without-RSU arms at every vehicle count. Message cost diverges
sharply and grows with load.

\begin{center}
\begin{tabular}{lrrrrrr}
\toprule
& \multicolumn{2}{c}{throughput (veh/s)} & \multicolumn{2}{c}{wait (s)}
& \multicolumn{2}{c}{msgs/vehicle} \\
\cmidrule(lr){2-3}\cmidrule(lr){4-5}\cmidrule(lr){6-7}
$N$ & OFF & ON & OFF & ON & OFF & ON \\
\midrule
4  & $0.557$ & $0.552$ & $5.64$  & $5.70$  & $65$  & $128$ \\
8  & $0.596$ & $0.595$ & $7.49$  & $7.49$  & $97$  & $181$ \\
12 & $0.677$ & $0.665$ & $9.65$  & $9.61$  & $159$ & $293$ \\
16 & $0.734$ & $0.723$ & $11.12$ & $11.26$ & $239$ & $374$ \\
20 & $0.741$ & $0.721$ & $13.24$ & $13.72$ & $319$ & $473$ \\
\bottomrule
\end{tabular}
\end{center}

\noindent
Message cost is $1.5$--$2.0\times$ higher with units at every count. Throughput
and wait differ by at most $0.02$ veh/s and $0.5$\,s respectively --- well
inside the repetition standard deviation, which is $0.04$--$0.18$ veh/s for
throughput and $0.8$--$1.8$\,s for wait.

Throughput \textbf{saturates above $N{\approx}16$} ($0.734 \to 0.741$ for OFF,
$0.723 \to 0.721$ for ON): the junction reaches capacity, and further vehicles
lengthen the queue without raising the clearing rate.

\subsection*{Why: units are not scheduled}

RSU units are non-vehicle PBFT replicas. They vote, echo and relay, but they
never queue at the junction. In \code{ResDBDecision.cc} the proposal packs unit
entries with \code{unit.is\_ambulance = 0} and they are never assigned a
crossing slot, so the \code{IntersectionExecutor} batches the same vehicles into
the same conflict-matrix groups in both arms.

The logs confirm this is not approximate but exact. At $N{=}16$ the departure
sequences are:

\begin{center}
\begin{tabular}{ll}
\toprule
OFF & 20.9, 27.2, 22.3, 28.6, 23.7, 21.9, \dots \\
ON  & 21.0, 27.3, 22.4, 28.7, 23.7, 22.0, \dots \\
\bottomrule
\end{tabular}
\end{center}

\noindent
The schedule is \textbf{byte-identical, uniformly displaced by $0.1$\,s}. That
offset is the extra consensus cost of four more replicas: PBFT ordering latency
rises from $0.213$\,s to $0.230$\,s at $N{=}16$, and the proposal is submitted
$0.178$\,s later. Against a ${\sim}10$\,s vehicle wait this is $1$--$2\%$.

\caveat{The top two panels cannot distinguish the arms by construction.} They
compare two executions of the same schedule and correctly find they take the
same time. No number of repetitions changes this --- the difference is
structural, not statistical.

\subsection*{Where the RSU benefit actually is}

Units raise quorum from $2f{+}1$ over $N$ vehicles to $2f{+}1$ over $N{+}4$
replicas --- verified in the logs as $3\to5$ at $N{=}4$ and $13\to15$ at
$N{=}20$. The measurable consequence is fault tolerance:

\begin{center}
\begin{tabular}{lll}
\toprule
& without RSU & with RSU \\
\midrule
$N{=}4$  & fails at $k{=}2$ & survives $k{=}2$ \\
$N{=}8$  & fails at $k{=}4$ & survives $k{=}4$ \\
$N{=}16$ & fails at $k{=}6$ & survives $k{=}6$ \\
$N{=}20$ & fails at $k{=}7$ & survives $k{=}7$; $2/3$ at $k{=}8$ \\
\bottomrule
\end{tabular}
\end{center}

\noindent
RSUs buy \textbf{one to two additional tolerated silent replicas} at every
vehicle count. \caveat{This is measured in 156 existing logs and is not
currently plotted.} Adding a commit-rate-versus-$k$ panel requires no new
simulation.

\caveat{Anomaly:} at $N{=}12$ the relationship inverts (OFF $k5{:}2/3$ versus ON
$k5{:}1/3$), the only count where units appear to hurt. Unexplained.

\subsection*{A defect this figure previously hid}

Before 27~August the RSU arm showed a large delay penalty (${\approx}18$\,s
versus $6$\,s at $N{=}4$). This was not a property of RSUs. Certificate
\emph{formation} derived $f$ from \code{total\_vehicles\_} while
\emph{validation} used \code{toleratedF()}, i.e.\ \code{num\_replicas\_}
including units. With units present these never agree, so every certificate was
assembled one to two echoes short of what every receiver demanded and every
replica dropped every other replica's certificate. Each vehicle retained only
its own, producing a proposal that was mostly QUIET, no commit, and a fall back
to stop-sign timeout. The arm was measuring the fallback path.

The defect was intermittent rather than constant because extra echoes sometimes
arrive before the certificate latches, which is why it presented as ``$N{=}4$
instability''. Three further sites used the same incorrect derivation and were
unified on \code{toleratedF()}: gossip propagation confirmation, consensus relay
carrier threshold, and a Byzantine-injection log line.

\section{Ablation 2 --- fake-ambulance attack}

\caveat{This study currently measures nothing and should not be presented.}

Colluders occupy the \emph{highest} replica ids, so \code{CertPrimary()}, which
elects the smallest certified id, always selects an honest proposer. This
isolates the certificate gate from leader succession. The arms are
\code{enableAmbulanceCertGate} on versus off.

\subsection*{Observation}

The gate fires correctly --- 264 rejections in the gate arm, 0 in the control.
But in the \emph{control} arm, where the certificate-less claim is deliberately
permitted, the liars gain nothing:

\begin{center}
\begin{tabular}{lll}
\toprule
& liar wait & honest wait \\
\midrule
no gate  & $12.86$\,s & $11.97$\,s \\
gate on  & $13.01$\,s & $12.16$\,s \\
\bottomrule
\end{tabular}
\end{center}

\noindent
The attackers wait ${\approx}0.9$\,s \emph{longer} than honest vehicles in both
arms, and both grant zero false priority.

\subsection*{Why}

The attack does not succeed even when undefended, so the gate is rejecting
something that was never exploitable. Byzantine type~7 sets
\code{isAmbulance = true} in the ARRIVAL\_ANNOUNCE with empty
\code{ambulanceCertBytes}. That claim is accepted into \code{local\_vehicle\_%
states\_}, but never becomes scheduling priority: \code{ResDBDecision.cc:146}
sets \code{e.is\_ambulance} from the proposer's vehicle states and the scheduler
prioritises on \code{is\_ambulance \&\& cyber\_status == 1}, yet no
\code{is\_ambulance=1} entry appears in any committed order. The path between
those two points drops the claim, and the cause has not been located.

Until that is traced, both arms will be flat regardless of the gate.

\section{Ablation 3 --- actuated traffic light versus consensus}

\textbf{9 repetitions.}

\subsection*{Observation}

\begin{center}
\begin{tabular}{lrrrr}
\toprule
$N$ & signal wait & ours & delay ratio & service-rate ratio \\
\midrule
8  & $12.61$\,s & $7.53$\,s & $1.67\times$ & --- \\
12 & $16.63$\,s & $9.69$\,s & $1.72\times$ & --- \\
16 & $16.85$\,s & $11.23$\,s & $1.50\times$ & --- \\
20 & $19.84$\,s & $13.20$\,s & $1.50\times$ & --- \\
\bottomrule
\end{tabular}
\end{center}

\noindent
The consensus arm is also \textbf{substantially more predictable}: standard
deviation $0.91$--$1.29$\,s versus $2.12$--$3.09$\,s for the signal, at every
vehicle count. The per-vehicle scatter panel shows the signal spanning
$4$--$44$\,s while the protocol stays within a band of roughly $10$--$20$\,s.

\subsection*{Why a signal, and why actuated}

Surveys of autonomous intersection management report fixed-time signalling as
the most-used baseline ($46\%$) while explicitly recommending actuated or
adaptive control instead; all-way-stop control is used by only $8\%$.

The all-way stop was additionally unusable here: under mixed turn demand it
deadlocked in 7 of 12 runs. Four vehicles at the head of each approach, all
turning left, each yield to the vehicle on their right, producing a circular
wait. With single-lane approaches this blocks following through-traffic as well.
Resolution came only from SUMO's $300$\,s teleport timer, producing
$300$--$600$\,s ``waits'' that are jam-resolution artifacts.

Fixed-time signalling would have been unfair \emph{in our favour}: with a single
arrival burst a fixed cycle burns green on approaches that have emptied and will
not refill. Actuated control extends green while vehicles are present and ends
the phase on a gap.

\subsection*{Why the protocol wins}

A signal grants right of way by \emph{phase}; the protocol computes a
\emph{schedule}. The conflict matrix admits simultaneous non-conflicting
movements that a two-phase signal must serialise across cycles, and there is no
inter-phase clearance interval. The variance difference has the same origin: a
signal's outcome depends on where in its cycle the burst arrives and how many
permissive left turns block its single-lane approaches, whereas a computed
schedule is insensitive to both.

\caveat{Protected left turns are impossible at this geometry.} Each approach has
one lane, so through and left movements share it and no phase can hold one red
while the other proceeds. \code{netconvert --tls.left-green.time} was verified to
produce a byte-identical signal program. This constrains the signal exactly as
it constrains the protocol.

\section{Ablation 4 --- ambulance priority}

\textbf{9 repetitions}, and the ambulance's route is \textbf{pinned}.

\subsection*{Why the methodology differs}

The ambulance is a \emph{single vehicle}, so this study samples $n{=}1$ per run
where the others average $8$--$20$. At 3 repetitions its wait varied by
$3.4$\,s between runs --- wider than the entire $N{=}8 \to 20$ trend. Its route
is additionally fixed to its approach's straight route while all other vehicles
remain random: a left-turning ambulance batches only with the opposing left,
whereas a straight one batches freely, so leaving the measured vehicle's turn to
chance let its own draw dominate the metric.

\subsection*{Observation}

\begin{center}
\begin{tabular}{lrrr}
\toprule
$N$ & priority & FIFO & benefit \\
\midrule
8  & $7.29$\,s & $10.34$\,s & $3.05$\,s \\
12 & $7.36$\,s & $13.93$\,s & $6.57$\,s \\
16 & $8.97$\,s & $15.72$\,s & $6.75$\,s \\
20 & $9.87$\,s & $18.52$\,s & $8.65$\,s \\
\bottomrule
\end{tabular}
\end{center}

\subsection*{Why the benefit grows with load}

The priority wait rises gently ($7.29 \to 9.87$\,s) while FIFO rises steeply
($10.34 \to 18.52$\,s). The ambulance is inserted near the front of the schedule
regardless of queue length, so what grows with $N$ is not its own wait but the
\emph{number of vehicles it skips}.

The cost is borne twice. Other vehicles wait ${\approx}2$\,s longer at $N{=}20$,
and the intersection's throughput falls from $0.75$ to $0.67$ vehicles/s ---
preempting the schedule breaks up batches that FIFO would have kept whole.

The \emph{relative} panel is the most defensible statement of the claim, since
dividing by the fleet mean in the same run cancels run-level variation.

\section{Ablation 5 --- late-emergency rollback}

\textbf{3 repetitions.} \caveat{Under-powered; see caveat below.}

\begin{center}
\begin{tabular}{lrrr}
\toprule
& cleared & wait & messages/vehicle \\
\midrule
rollback OFF & $16.0$ & $10.40$\,s & $357$ \\
rollback ON  & $16.0$ & $12.30$\,s & $554$ \\
\bottomrule
\end{tabular}
\end{center}

\subsection*{Why both arms clear the same number}

Rollback does not serve \emph{fewer} vehicles; it serves \emph{different} ones.
The ambulance occupies a slot an ordinary vehicle would have taken and the
remainder shift later. Every count-based summary therefore shows the arms as
identical, which is why the left panel plots individual departures rather than a
total --- the substitution is invisible to any statistic over counts.

The cost is ${\approx}1.9$\,s of additional wait and $1.55\times$ the messages:
revising an already-committed order requires a CANCEL round on top of the ORDER
round.

\caveat{Rollback fired in 2 of 3 repetitions in the latest run and 3 of 3
previously.} At 3 repetitions a regression cannot be distinguished from the turn
draw. This is the cheapest study (${\approx}5$\,min per 3 repetitions) and
should be run at 9--12.

\section{Cost decomposition}

Honest operation only ($k{=}0$): with replicas silenced the traffic mix reflects
failure rather than normal composition.

\subsection*{Observation: the two panels point in opposite directions}

Per occurrence, \textbf{PBFT ordering overtakes certificates in message count}
as $N$ grows --- a small fraction at $N{=}4$, over half the stack at $N{=}20$.
This follows from the structure: one PBFT round costs $2n{+}1$ sends
(PRE\_PREPARE broadcast, then PREPARE and COMMIT from each replica) while one
certificate costs $n{+}2$ (one announce, $f{+}1$ echoes, one certificate
broadcast).

Yet \textbf{certificate formation takes longer} --- $0.30$\,s versus $0.24$\,s
at $N{=}20$.

\subsection*{Why: certificate latency is bounded by the announce period}

A replica emits an ARRIVAL\_ECHO only when it \emph{hears} an announce; it has
no other trigger. A vehicle therefore cannot gather its $(f{+}1)$th echo until it
has announced often enough for that many distinct peers to have heard it. Peers
that come into range late, or whose earlier echo was lost, must wait for the
\emph{next} announce.

Announce transmission is exponentially backed off ($0.1$\,s doubling to a
$3.2$\,s ceiling, reset when a previously unseen vehicle is registered), so that
wait is bounded by the current backoff interval rather than by network latency.
The correlation in a single $N{=}16$ run is close to linear:

\begin{center}
\begin{tabular}{rr}
\toprule
announce interval & certificate latency \\
\midrule
$0.20$\,s & $0.028$\,s \\
$0.25$\,s & $0.033$\,s \\
$0.40$\,s & $0.342$\,s \\
$0.80$\,s & $0.694$\,s \\
\bottomrule
\end{tabular}
\end{center}

\noindent
Certificate formation is therefore \emph{not} transmission-bound. It is
dominated by waiting for the vehicle's own next scheduled announce.

\subsection*{The backoff trade, measured}

A controlled A/B (identical binary, seed and routes; only the backoff ceiling
differs) at $N{=}16$:

\begin{center}
\begin{tabular}{lrrr}
\toprule
& backoff on & backoff off & change \\
\midrule
certificate-layer messages & $3\,132$ & $10\,764$ & $-71\%$ \\
certificate-layer bytes    & $1.55$\,M & $2.95$\,M & $-47\%$ \\
all messages               & $3\,737$ & $11\,422$ & $-67\%$ \\
certificate latency (median) & $0.091$\,s & $0.046$\,s & $2.0\times$ worse \\
\textbf{vehicle wait}      & $9.81$\,s & $9.89$\,s & $-0.8\%$ \\
\bottomrule
\end{tabular}
\end{center}

\noindent
The backoff removes two-thirds of all radio traffic at a latency cost that
\textbf{does not reach the vehicle}: end-to-end wait differs by $0.8\%$, in the
favourable direction. Certificate formation is ${\sim}0.1$\,s against a
${\sim}10$\,s journey, three orders of magnitude below the quantity that
matters. The trade is strongly favourable, and more so on a shared channel with
real interference than in this clean simulation.

\caveat{Scope:} one seed at $N{=}16$. The latency distribution is
right-skewed (mean $0.129$\,s, median $0.091$\,s), so the median is the honest
summary.

\section{Summary of open issues}

\begin{enumerate}
\item \textbf{Ablation 2 measures nothing} --- the attack does not succeed even
      undefended; the announce-to-scheduler path must be traced.
\item \textbf{Ablation 1 does not plot the fault frontier}, which is the sole
      panel able to distinguish the arms. Data already exists.
\item \textbf{Ablation 1, $N{=}12$}: fault tolerance inverts with units.
\item \textbf{Ablation 5} requires more repetitions --- rollback fired in 2 of 3
      runs in the latest execution and 3 of 3 previously, which 3 repetitions
      cannot separate from the turn draw.
\item \textbf{Single-lane approaches} preclude protected-left phases for all
      arms; results do not generalise to a junction with turn pockets.
\item \textbf{Burst arrival, single epoch}: no steady-state comparison is
      possible until multi-epoch support exists.
\end{enumerate}

\end{document}
